\section{Removing the comparison-output restriction}
\label{sec:risk-weighted-multiplier}
A variance-sensitive multiplier estimate removes the output cap from the
preceding arbitrary-input-law oracle without worsening its rate.
The fixed comparison is selected using population risk, before sampling;
it need not have a uniformly small output.

Keep the Gaussian model, activation, resource $A_n(g)$ in (AD1), and
input cube of the preceding section, but now allow any finite
comparison network $g$, with no output restriction. Write
$V=\sum_i|c_i|$ and $F=\|g\|_\infty\le K V$.
For the proof it suffices to consider $R(g)\le R(0)+1\le2$ under
bounded labels: a comparison with larger risk can be replaced in the
proof by the zero network, at the same likelihood scale.

\begin{lemma}[Risk-weighted ridge multipliers]
\label{lem:risk-weighted-ridge}
Let $u$ be a fixed measurable function with $\|u\|_\infty\le F_u$,
let $|Y|\le B$, and put
$m(X,Y)=u(X)-Y$ and $s^2=\mathbb E m^2$.
For $t=\log(2/\delta)$, with probability at least $1-\delta$,
\[
 \sup_\theta\bigl|(P-P_n)[m\phi(\theta\cdot\widetilde X)]\bigr|
 \le C_\phi\left\{
       s\sqrt{(D+t)/n}+(F_u+B)t/n\right\}=:\eta_n.       \tag{WM1}
\]
The inequality holds simultaneously over every incoming weight vector,
without a norm bound or any independence between $m$ and $X$.
\end{lemma}
\begin{proof}
Condition on the observed pairs and write $m_i=m(X_i,Y_i)$.
The weighted Rademacher process has increment metric
\[
 d_m(h,h')^2=n^{-2}\sum_i m_i^2(h(X_i)-h'(X_i))^2.
\]
If $\sum_i m_i^2>0$, normalize these weights to a probability measure
$P_{n,m}=\sum_i m_i^2\delta_{X_i}/\sum_i m_i^2$.
The VC-subgraph covering bound for the bounded ridge class is valid for
this probability measure too. Dudley's integral therefore gives
\[
 \mathbb E_\epsilon\sup_\theta
 \left|n^{-1}\sum_i\epsilon_i m_i\phi(\theta\cdot\widetilde X_i)\right|
 \le C_\phi\sqrt D\,\frac{\sqrt{\sum_i m_i^2}}n.
\]
For zero total weight the process is identically zero. Taking expectation,
using Jensen, and symmetrizing shows that the expected centered supremum
is at most $C_\phi s\sqrt{D/n}$.

For the class with both signs, each summand has envelope
$U=K(F_u+B)$ and variance at most $K^2s^2$.
Bousquet's inequality~\cite[Theorem~2.3]{bousquet} bounds its centered supremum by its expectation
plus
\[
 C\sqrt{(K^2s^2+U\mathbb E Z)t/n}+CUt/n.
\]
The mixed term $\sqrt{U(\mathbb E Z)t/n}$ is at most
$C\mathbb E Z+CUt/n$, by Young's inequality.
This proves (WM1). The weighted covering argument is conditional on the
observed pairs and uses only $m_i^2$, so dependence on $X$ causes no issue.
\end{proof}

\begin{theorem}[Sharp uncapped-comparison oracle]
\label{thm:uncapped-all-d}
Assume $|Y|\le1$, and fix any comparison network $g$.
Let $b_n$ be (AD4), with a sufficiently large
activation-dependent constant. In the nontrivial range
$A_n(g)b_n\le c_\phi$, with probability at least $1-\delta$,
simultaneously for all $0<\nu\le1$, all
$n^{-2}\le\lambda\le1$ with $\lambda\ge C_\phi b_n$, and all
$N\lambda\ge C_\phi$,
\[
 R(\mathbb E_{\Pi_\lambda}f)\le\Pi_\lambda R(f)
 \le R(g)+C_\phi\{\lambda A_n(g)+b_n/\lambda\}
       +C_\phi n^{-10}.                                \tag{WM2}
\]
In particular, $\lambda=C_\phi\sqrt{b_n/A_n(g)}$ yields
\[
 R(\mathbb E_{\Pi_\lambda}f)\le\Pi_\lambda R(f)
 \le R(g)+\widetilde O_\phi\left(\sqrt{(d+1)A_n(g)/n}\right),         \tag{WM3}
\]
with the same learner-width condition as (AD11). There is no output
cap on either the comparison network or the learner.
\end{theorem}
\begin{proof}
It suffices to prove the bound when $R(g)\le R(0)+1$. Otherwise
apply that case to zero at the same $\lambda$: $A_n(0)\le A_n(g)$
and $R(0)\le R(g)$, so its conclusion is stronger than (WM2).
This is a proof choice and does not require the learner to know either risk.

Use the variation-weighted comparison pack from (AD8), with
$\sigma=n^{-20}/(V+1)$, and write $u=g_\sigma$.
Its uniform error is $O_\phi(n^{-20})$. By the triangle inequality in
population $L^2$, its residual risk satisfies
\[
 \sqrt{R(u)}\le\sqrt{R(g)}+C_\phi n^{-20}.
\]
Thus $R(u)$ is bounded by an activation-dependent constant,
without using the output envelope $F$. Moreover $u$ is a signed
Gaussian mixture of ridge features with total variation exactly $V$
(or at most $V$ after combining coincident features).

Apply (WM1) to $m=u-Y$. Its event is uniform over every ridge, so for
any learner function $f\in\mathcal F_q$,
\[
 |(P-P_n)[m(f-u)]|\le(\sqrt q+V)\eta_n.                 \tag{WM4}
\]
For the $u$ term this follows by integrating the common ridge bound
against its signed mixing measure; boundedness justifies Fubini.
For $f$ it follows by its finite signed-hull representation.
Since $\|u\|_\infty\le F+C_\phi n^{-20}$,
\[
 \eta_n^2\le C_\phi\{b_n+(F+1)^2b_n^2\}
 \le C_\phi b_n.                                      \tag{WM5}
\]
The last step uses $(F+1)^2\le C_\phi A_n$ and $A_nb_n\le c_\phi$.
Young's inequality in (WM4), with a chosen small constant, gives
\[
 2|(P-P_n)[m(f-u)]|
 \le c\lambda q+C\lambda V^2+C_\phi b_n/\lambda.        \tag{WM6}
\]

The smooth-loss norm comparison (AD6) does not need $u$ to be bounded
by two. With its actual envelope it gives
\[
 (P_X-P_n)(f-u)^2
 \le\tfrac16P_n(f-u)^2
       +C_\phi b_n\{q+(F+1)^2\}.                      \tag{WM7}
\]
Its $q$ term is absorbed into a small multiple of $\lambda q$ by
$\lambda\ge C_\phi b_n$; its constant term is at most
$C_\phi\lambda A_n$. Combine (WM6)--(WM7) with the relative-risk
identity used in (AD5), integrate, and use the exact midpoint
comparison (AD9). That comparison and its pack entropy do not require
an output cap. Its empirical-distance and readout-energy penalties
absorb the chosen small terms, proving (WM2).

For completeness, the same finite energy cutoff is valid here.
The condition $A_nb_n\le c$ implies $F^2\le C_\phi A_n\lesssim n$.
The discarded distance and relative-risk integrands are bounded by
$C_\phi(1+Q_G+F^2)$, so the original Gaussian readout tail above
$n^{22}$ remains negligible uniformly over the claimed scales.
Finally $R(u)\le R(g)+C_\phi n^{-20}$ by the residual $L^2$ inequality,
and Jensen gives the raw-mean claim. All sample events use only the
fixed $u$ and shared feature classes, preserving simultaneity in $N,\nu$.
\end{proof}

\paragraph{Marginally sub-Gaussian labels and the class infimum.}
Suppose instead $\mathbb E\exp(Y^2/K_Y^2)\le2$.
Fix any uncapped comparison network. As above, comparisons with
$R(g)>R(0)+1$ are replaced by zero only in the proof.
The same conclusion holds with
\[
 b_n^*=C_{\phi,K_Y}(D+1)\log^5(e n/\delta)/n
\]
in place of $b_n$, under the same nontrivial-range and width conditions.
To verify it, use the proof-only label cutoff from (AD14),
$B_n=\max\{1,K_Y\sqrt{40\log(e n/\delta)}\}$.
The clipped-label residual has bounded second moment because
\[
 \|g_\sigma-Y^{B_n}\|_{L^2}
 \le\sqrt{R(g)}+\|Y-Y^{B_n}\|_{L^2}+C_\phi n^{-20}
 \le C_{\phi,K_Y}.
\]
Apply (WM1) with envelope $F+C_\phi+B_n$. Its squared bound is
$C_{\phi,K_Y}\{b_n^*+(F+1)^2(b_n^*)^2\}$, absorbed by
$C_{\phi,K_Y}b_n^*$ when $A_nb_n^*$ is small.
The other proof steps are unchanged. The actual posterior agrees with
the proof's clipped-label posterior on the no-exceedance event, and
\[
 [R_Y(f)-R_Y(g)]-[R_{Y^{B_n}}(f)-R_{Y^{B_n}}(g)]
 =-2\mathbb E[(Y-Y^{B_n})(f(X)-g(X))].
\]
Here $F\lesssim\sqrt{A_n}\lesssim\sqrt n$, and the readout energy
moment from (AD14) bounds the posterior average of $\|f\|_\infty$.
Thus the marginal tail still makes this transfer $O_{\phi,K_Y}(n^{-10})$.
No conditional-mean assumption and no actual label clipping are used.

For any declared coefficient-bounded shallow teacher class containing
zero, select a deterministic population near-minimizer; it satisfies
$R(g)\le R(0)+1$. Equations (WM2)--(WM3) therefore hold with its
risk replaced by the class infimum and a deterministic upper bound on
$A_n$, with exact leading coefficient one. For fixed $d,B$, the error
is $\widetilde O(C/\sqrt n)$ against \emph{all} width-$C$ teachers
with coefficient cap $B$, not just those with bounded output.
In particular, the single schedule $\lambda=n^{-3/4}$ competes with
all width-$n^{1/4}$ such teachers at excess
$\widetilde O(n^{-1/4})$, simultaneously for every $N\ge c_\phi n$
and $0<\nu\le1$, once $n$ is sufficiently large. The schedule does not
use the unknown comparator coefficients or risk.
The same removal of the output cap applies to the exact ReLU-to-hardtanh
compiler in the preceding section.

\begin{corollary}[Growing width and coefficient budgets]
\label{cor:all-d-growing-resources}
Fix $d$, the activation, a marginal sub-Gaussian label envelope, and
nonnegative exponents $a,b$ with $a+b<1/2$. For a class of shallow
teachers of width at most $n^a$ and coefficient cap $n^b$, with no
output cap, a deterministic scale
$\lambda_n=\widetilde\Theta(n^{-1/2-a-b})$ gives
\[
 R(\mathbb E_{\Pi_{\lambda_n}}f)\le\Pi_{\lambda_n}R(f)
 \le\inf_g R(g)+\widetilde O(n^{a+b-1/2})
\]
with probability at least $1-\delta$, simultaneously over all
$0<\nu\le1$ and $N\ge c_\phi n$, for sufficiently large $n$.
The scale uses the declared growth budgets, not the unknown teacher.
For a hardtanh learner against coefficient-capped ReLU teachers, the
corresponding condition is $a+2b<1/2$, with remainder
$\widetilde O(n^{a+2b-1/2})$.
\end{corollary}
\begin{proof}
Equation (AD2) gives $A_n=\widetilde O(n^{2(a+b)})$ at fixed $d$.
Balance the two terms of (WM2). Its nontrivial-range and width conditions
hold because $a+b<1/2$. The hardtanh compiler instead has
$V'=O_d(n^{a+2b})$ and $J'=O_d(1)$, giving the stated substitution.
\end{proof}

There is also one schedule that needs neither growth exponent. At fixed
$d$, put $\ell_n=\log(e n/\delta)$ and
$\lambda_n=\ell_n^7/n$. For sufficiently large $n$, (WM2) gives
\begin{equation}
 R(\mathbb E_{\Pi_{\lambda_n}}f)\le\Pi_{\lambda_n}R(f)
 \le\inf_{g\in\mathcal G_n}R(g)
 +C_{\phi,d,K_Y}\left\{\frac{a_n\ell_n^7}{n}+\ell_n^{-2}\right\}
 +C_{\phi,d,K_Y}n^{-10},                              \tag{WM8}
\end{equation}
for any declared class $\mathcal G_n$ containing zero with
$\sup_{g\in\mathcal G_n}A_n(g)\le a_n$ and
$a_n b_n^*\le c_\phi$. The event is common to every eligible learner
width $N\ge c_\phi n$ and scale $0<\nu\le1$; the algorithm does not
use $a_n$. Indeed $b_n^*/\lambda_n=O_{\phi,d,K_Y}(\ell_n^{-2})$,
while the first term is $\lambda_n a_n$. Thus this same schedule gives
a vanishing remainder for each fixed pair $a,b$ with $a+b<1/2$ above.
This means one learning rule has the stated guarantee for each declared
class, not that all data-selected comparison networks share a single
unproved event.
